\scanpage{113}
% archive foliation: 112
ou encore, en posant $\bar\omega=\lambda(t)\,dt$
\[
\frac{s''}{s'}=-\lambda(t)\qquad\text{i.e.}\quad(\log s')'=-\lambda(t)
\]
i.e. $\log s'=-\Lambda(t)+c^{\mathrm{te}}$ ($\Lambda(t)$ une primitive de $\lambda(t)$)
\[
s'=\underset{\ne0}{c^{\mathrm{te}}}\exp\Lambda(t)\qquad
s=c^{\mathrm{te}}\int\exp-\Lambda(t)\,dt+c^{\mathrm{te}}.
\]
donc les solutions existent (analytiques) et \deleted{sont} forment un torseur \inserted{analytique} sous le gr. affine de $\operatorname{Aut}(\mathbb E^1(\mathbb C))=(\operatorname{Aut}\mathbb E^1_{\mathbb C})^{\mathrm{an}}$ ; c'est le \deleted{\unsure{du}} l'interprétation transcendante du torseur \deleted{$\unsure{\longrightarrow}$} sous $B_{\mathbb C}^{\mathrm{an}\deleted{\unsure{-1}}}$ \deleted{de} d'image inverse dans $\mathcal T$ de la section donnée $\sigma'$ de $P$. Le torseur quotient de celui-ci par $B_u^{\mathrm{an}}\simeq\mathbb G_a^{\mathrm{an}}$ est le suivant :
\inserted{une section en dehors de $\unsure{\sigma'|\sigma}$ définie par $D$ est}
\deleted{on prend les} faisceau de \unsure{l'ens.} des 1-formes $\varphi$ de la forme $ds$, où $D_s=\underline D$, i.e. des formes $\varphi$ telles que $\uline{D\varphi=\deleted{\unsure{\varphi}}0}$, $\uline{\varphi\text{ ne s'annule pas}}$.
\margin{?}

\margin{! \unsure{description alg.} \deleted{$\unsure{\triangle}$}}
\deleted{alg.} Sections $\sigma'$ \inserted{\unsure{algébriques}} de $P\inserted{=\mathcal T/B}$ horizontales disjointes de $\sigma$ : connexions sur $\underline\omega$ telles que $c_D=c_0$.

\margin{description transcendante}
Sections \uline{analytiques} de $\mathcal T/B_u$ horizontales, \unsure{associées} à sections $\sigma'$ disjointes de $\sigma$ : couples \deleted{de} $(D,\varphi)$, $D$ connexion \inserted{\unsure{analyt. sur $\underline\omega$}} telle que $c_D=c_0$, $\varphi\in\Gamma^{\mathrm{an}}\underline\omega$ telles que $D\varphi=0$.

Sections analytiques de $\mathcal T$ horizontales \inserted{\unsure{telles que $\sigma'$ disj. de $\sigma$}} : uniformisants
\[
s:\deleted{\unsure{\mathbb E}}X^{\mathrm{an}}\longrightarrow\mathbb E^1_{\mathbb C}
\quad\text{\deleted{tels que} tels que}\quad s^*(c)=c_0
\]
i.e. $c_{D_t}=c_0$
\[
t\longmapsto(D_t,dt)\longmapsto D_t
\]
correspond à les applications canoniques
\diagram{\begin{tikzcd}
\mathcal T\arrow[r]&\mathcal T/B_u\arrow[r]&\mathcal T/B
\end{tikzcd}}
